\documentclass[12pt]{letter}
\usepackage[margin=1in]{geometry}

\begin{document}

\begin{letter}{Editor-in-Chief\\
Knowledge and Information Systems\\
Springer Nature}

\opening{Dear Editor,}

Please consider my manuscript, entitled \textbf{``Retrieval-Assisted Instantiation of Natural-Language Optimization Problems,''} for publication in \textit{Knowledge and Information Systems}.

The paper studies a knowledge-processing problem that sits between natural-language understanding and formal optimization modeling: given a natural-language description of an optimization problem, how much structured, usable knowledge -- a compatible optimization schema and its numeric parameters -- can be recovered by a transparent, deterministic pipeline, and where does that recovery break down? Rather than targeting full end-to-end solver-ready code generation, the manuscript isolates two intermediate capabilities that are often conflated in end-to-end systems: (i) retrieval of a compatible optimization schema from a fixed catalog, and (ii) schema-conditioned, typed grounding of numeric evidence from text into that schema's scalar parameters.

On the \textsc{NLP4LP} benchmark, deterministic lexical retrieval (TF--IDF, BM25, LSA) already identifies the correct schema in the great majority of cases (TF--IDF Schema R@1 $=0.9094$), while an oracle control with perfect retrieval improves the end-to-end \textit{InstantiationReady} score only modestly, from $0.5287$ to $0.5680$. This gap is the central empirical finding of the paper: once a plausible schema is available, semantically grounding extracted numbers into the correct parameter roles -- not schema selection -- is the dominant bottleneck. The manuscript supports this claim with a numeric-type-compatibility ablation of the numeric type-handling logic, statistical significance testing, lexical-overlap stress tests, an error taxonomy, and three restricted engineering-oriented validation studies (a 60-instance structural subset, a 269-instance executable-attempt subset that transparently documents an environment blocker, and a 20-instance solver-backed subset using a Gurobi-free SciPy HiGHS compatibility shim).

I believe the manuscript fits the scope of \textit{Knowledge and Information Systems} because it addresses structured knowledge representation and retrieval-assisted semantic grounding for optimization-relevant natural-language content -- an intelligent information-processing problem -- and because it reports a controlled, reproducible benchmark study with an honest bottleneck analysis rather than an end-to-end automation claim. The contribution is explicitly scoped: a deterministic retrieval-and-grounding framework, a retrieval-dependent evaluation design that separates schema-selection error from grounding error, and a transparent, heuristic-but-well-specified compatibility procedure for scalar parameter instantiation.

This submission is original, has not been published previously, and is not under consideration for publication elsewhere. All benchmark values reported in the manuscript are drawn from measured artifacts in the accompanying code repository, with provenance documented in the repository's analysis reports.

Thank you for your time and consideration. I appreciate the opportunity to submit this work to \textit{Knowledge and Information Systems}.

\closing{Sincerely,\\[1em]
Soroush Vahidi\\
Ying Wu College of Computing\\
New Jersey Institute of Technology\\
Newark, NJ, USA\\
Email: sv96@njit.edu\\
ORCID: 0000-0003-1934-6282}

\end{letter}

\end{document}
